\usepackage[T1]{fontenc}
\usepackage[utf8]{inputenc}
\usepackage{lmodern}
\usepackage[margin=25mm]{geometry}
\usepackage{amsmath,amssymb,amsthm}
\usepackage{graphicx,xcolor,enumitem,booktabs,array}
\usepackage{microtype,titlesec}
\usepackage[breakable,skins]{tcolorbox}
\usepackage[hidelinks,unicode]{hyperref}
\definecolor{notesblue}{HTML}{335C81}
\definecolor{notesamber}{HTML}{9A6926}
\definecolor{notesink}{HTML}{253746}
\titleformat{\section}{\Large\bfseries\raggedright\color{notesblue}}{\thesection}{0.7em}{}
\titleformat{\subsection}{\large\bfseries\raggedright\color{notesink}}{\thesubsection}{0.7em}{}
\titleformat{\subsubsection}{\normalsize\bfseries\raggedright\color{notesink}}{\thesubsubsection}{0.7em}{}
\titleformat{\paragraph}[block]{\normalsize\bfseries\raggedright\color{notesink}}{\theparagraph}{0.7em}{}
\titlespacing*{\section}{0pt}{18pt plus 3pt minus 2pt}{8pt}
\titlespacing*{\subsection}{0pt}{12pt plus 2pt minus 2pt}{6pt}
\titlespacing*{\subsubsection}{0pt}{9pt plus 2pt minus 1pt}{4pt}
\titlespacing*{\paragraph}{0pt}{9pt plus 2pt minus 1pt}{4pt}
\setlength{\parindent}{0pt}
\setlength{\parskip}{6pt plus 2pt minus 1pt}
\setlength{\emergencystretch}{3em}
\setlist{leftmargin=*,topsep=4pt,itemsep=2pt,parsep=2pt,partopsep=0pt}
\renewcommand{\arraystretch}{1.15}
\AtBeginDocument{%
  \setlength{\abovedisplayskip}{8pt plus 2pt minus 2pt}%
  \setlength{\belowdisplayskip}{8pt plus 2pt minus 2pt}%
  \setlength{\abovedisplayshortskip}{4pt plus 2pt}%
  \setlength{\belowdisplayshortskip}{6pt plus 2pt minus 2pt}%
}
% Unnumbered presentation wrappers: labels and statements remain in the source.
\tcbset{notesbox/.style={enhanced,breakable,sharp corners,
  boxrule=0.4pt,leftrule=1.6pt,colframe=notesblue!65,
  colback=notesblue!3,coltext=black,
  left=9pt,right=9pt,top=7pt,bottom=7pt,boxsep=0pt,
  before skip=9pt plus 2pt,after skip=9pt plus 2pt,
  lines before break=20,pad at break*=3pt}}
\newtcolorbox{notesdefinition}{notesbox}
\newtcolorbox{notestheorem}{notesbox,colback=notesblue!5,colframe=notesblue}
\newtcolorbox{notesresult}{notesbox,colback=notesblue!5,colframe=notesblue}
% The shared break threshold keeps short statements and algorithms together.
% Longer blocks can still span output pages.
\newtcolorbox{notesalgorithm}{notesbox,colback=notesamber!7,colframe=notesamber!80}
% Problem statements use the source PDFs' large left brace and math alignment.
\newenvironment{notesproblem}{%
  \[\left\{\begin{aligned}%
}{%
  \end{aligned}\right.\]%
}
\newlength{\algorithmindent}
\setlength{\algorithmindent}{1.5em}
\newenvironment{notespseudocode}{%
  \par\begingroup\setlength{\parskip}{2pt}%
}{\par\endgroup}
% Each source line can wrap, with continuation text offset from its first line.
\newcommand{\pseudocodeline}[2]{%
  \par\begingroup
  \leftskip=\dimexpr#1\algorithmindent+\algorithmindent\relax
  \hangindent=\algorithmindent\hangafter=1
  \noindent\strut#2\par\endgroup
}
\newcommand{\sourcepage}[2]{\par\phantomsection\label{source:#1:#2}}
\newcommand{\sourceheading}[2][]{%
  \subsection*{#2}%
  \if\relax\detokenize{#1}\relax
    \addcontentsline{toc}{subsection}{#2}%
  \else
    \addcontentsline{toc}{subsection}{#1}%
  \fi
}
\newcommand{\lecturecontents}[1]{\phantomsection\addcontentsline{toc}{part}{#1}}
\newcommand{\sourcesubheading}[1]{\paragraph*{#1}}
\newcommand{\sourcefigure}[2][0.9]{\par\begin{center}\includegraphics[width=#1\linewidth,height=0.85\textheight,keepaspectratio]{#2}\end{center}\par}
\newcommand{\bits}{\{0,1\}}
\newcommand{\eps}{\varepsilon}
\newcommand{\polyred}{\leq_p}
\newcommand{\classP}{\mathrm{P}}
\newcommand{\classNP}{\mathrm{NP}}
\newcommand{\classNPC}{\mathrm{NPC}}
\newcommand{\classNPH}{\mathrm{NPH}}
\newcommand{\problem}[1]{\textnormal{\textsc{#1}}}
\newcommand{\code}[1]{\texttt{#1}}
\newcommand{\xRightarrow}[1]{\overset{#1}{\Longrightarrow}}
\raggedbottom
\newcommand{\sourcecontinuation}[2]{\phantomsection\label{source:#1:#2}}
